% 全国大学生统计建模大赛 paper skeleton for a Mathodology workspace.
% Rules: 2026 (12th) 参赛须知 for the provincial round, read from a university copy
% on 2026-11-05 (the official site could not be read):
%   https://cxcyxy.ntu.edu.cn/_upload/article/files/a8/d9/3d3c581042f5ab1a36794cac860c/9b42dbdd-eb89-4874-993a-7b2f00329ea5.pdf
% Not an official template. The contest also wants a Word full version; this file
% gives the PDF. See ../../../../references/contest-formats.md. Compile with XeLaTeX.
%
% \anonymoustrue  -> PDF anonymous version: no cover page, no acknowledgements.
% School, team and advisor may appear only on the cover and in the acknowledgements.
% Contents to three levels plus a list of tables and figures. Body pages are
% numbered separately from 1. 3-5 keywords.
\documentclass[UTF8,a4paper,zihao=-4]{ctexart}
\usepackage[top=2.54cm,bottom=2.54cm,left=3.17cm,right=3.17cm]{geometry}
\usepackage{amsmath,amssymb,booktabs,graphicx}
\usepackage[hidelinks]{hyperref}
\usepackage{gbt7714}
\bibliographystyle{gbt7714-numerical}
\graphicspath{{../figures/}}
\input{results}

\newif\ifanonymous \anonymoustrue
\newif\ifenglishabstract \englishabstractfalse   % not required by the 2026 notice

\newcommand{\ContestYear}{2027}
\newcommand{\Edition}{十三}
\newcommand{\WorkNumber}{作品编号}
\newcommand{\PaperTitle}{论文题目}
\newcommand{\School}{参赛学校}
\newcommand{\Members}{队员一\quad 队员二\quad 队员三}
\newcommand{\Advisor}{指导老师}

% 2026 notice: level 1 黑体小三, level 2 楷体四号, level 3 宋体小四加粗, body 宋体小四,
% fixed 24pt leading (approximated here; check the compiled pages).
\linespread{1.55}
\ctexset{
  section={format=\heiti\zihao{-3},name={,、},number=\chinese{section},aftername={}},
  subsection={format=\kaishu\zihao{4},name={（,）},number=\chinese{subsection},aftername={}},
  subsubsection={format=\songti\zihao{-4}\bfseries,number=\arabic{subsubsection}.,aftername={}},
}
\setcounter{tocdepth}{3}
\pagestyle{plain}

\begin{document}

\ifanonymous\else
\thispagestyle{empty}
{\heiti\zihao{3} 作品编号：\WorkNumber}\par
\vspace*{5cm}
\begin{center}
{\heiti\zihao{-2} \ContestYear 年（第\Edition 届）全国大学生统计建模大赛}\par\vspace{20pt}
{\zihao{1} 参\ 赛\ 作\ 品}\par\vspace{4cm}
\zihao{-3}\renewcommand{\arraystretch}{2}
\begin{tabular}{rl}
参赛学校： & \School \\
论文题目： & \PaperTitle \\
参赛队员： & \Members \\
指导老师： & \Advisor \\
\end{tabular}
\end{center}
\clearpage
\fi

\pagenumbering{Roman}
\begin{center}{\heiti\zihao{3} \PaperTitle}\end{center}
\section*{摘\quad 要}
\addcontentsline{toc}{section}{摘要}
研究目的、方法、结果、结论与创新点。

\noindent{\heiti 关键词：}关键词一；关键词二；关键词三

\ifenglishabstract
\clearpage
\section*{Abstract}
\addcontentsline{toc}{section}{Abstract}
\fi

\clearpage
\tableofcontents
\clearpage
\section*{表格与插图清单}
\addcontentsline{toc}{section}{表格与插图清单}
\listoftables
\listoffigures

\clearpage
\pagenumbering{arabic}
\begin{center}{\heiti\zihao{3} \PaperTitle}\end{center}
\section{引言}
\section{数据来源与处理}
\section{研究方法}
\section{实证结果}
\section{结论与建议}

% \bibliography{references}

\clearpage
\section*{附录}
\addcontentsline{toc}{section}{附录}

\ifanonymous\else
\clearpage
\section*{致谢}
\addcontentsline{toc}{section}{致谢}
\fi
\end{document}
